%%%%%%%%%%%%%%%%%%%%%%%%%%%%%%%%%%%%%%%%%
% Jacobs Portrait Poster
% LaTeX Template
% Version 1.0 (31/08/2015)
%%%%%%%%%%%%%%%%%%%%%%%%%%%%%%%%%%%%%%%%%

\documentclass[final]{beamer}

\usepackage[scale=1.24]{beamerposter}
\usetheme{confposter}

\setbeamercolor{block title}{fg=ngreen,bg=white}
\setbeamercolor{block body}{fg=black,bg=white}
\setbeamercolor{block alerted title}{fg=white,bg=dblue!70}
\setbeamercolor{block alerted body}{fg=black,bg=dblue!10}

%-----------------------------------------------------------
% Dimensiones del póster
%-----------------------------------------------------------
\newlength{\sepwid}
\newlength{\onecolwid}
\newlength{\twocolwid}
\newlength{\threecolwid}
\setlength{\paperwidth}{36in}
\setlength{\paperheight}{48in}
\setlength{\sepwid}{0.024\paperwidth}
\setlength{\onecolwid}{0.22\paperwidth}
\setlength{\twocolwid}{0.464\paperwidth}
\setlength{\threecolwid}{0.708\paperwidth}
\setlength{\topmargin}{-0.5in}
%-----------------------------------------------------------

\usepackage[utf8]{inputenc}
\usepackage[T1]{fontenc}
\usepackage[spanish]{babel}
\usepackage{graphicx}
\usepackage{booktabs}
\usepackage{xcolor}
\usepackage{url}
\graphicspath{ {./images/} }

%-----------------------------------------------------------
% Marcador de datos faltantes.
% Buscar "PENDIENTE:" con Ctrl+F para completar uno por uno.
%-----------------------------------------------------------
\newcommand{\pendiente}[1]{\textcolor{red}{\textbf{[PENDIENTE: #1]}}}

%----------------------------------------------------------------------------------------
% TÍTULO
%----------------------------------------------------------------------------------------
\title{Sistema de Comunicación Morse mediante Audio}
\author{Cristian Juárez, Nicolaus Hernán y Zambrano Serrano Héctor}
\institute{Universidad Nacional de Hurlingham \\ Arquitectura de Computadoras --- Docente: Ing. Caroli Natalia}

%----------------------------------------------------------------------------------------

\begin{document}

\addtobeamertemplate{block end}{}{\vspace*{2ex}}
\addtobeamertemplate{block alerted end}{}{\vspace*{2ex}}
\setlength{\belowcaptionskip}{2ex}
\setlength\belowdisplayshortskip{2ex}

\begin{frame}[t]
\begin{columns}[t]

\begin{column}{\sepwid}\end{column}

%============================================================
% COLUMNA 1
%============================================================
\begin{column}{\onecolwid}

%----------------------------------------------------------------------------------------
% OBJETIVOS
%----------------------------------------------------------------------------------------
\begin{alertblock}{Objetivos}
\begin{itemize}
\item Recibir señales acústicas correspondientes a código Morse emitidas por otro dispositivo.
\item Detectar la frecuencia del tono recibido mediante el algoritmo de Goertzel.
\item Diferenciar puntos, rayas y pausas a partir de la duración de los pulsos.
\item Decodificar el código Morse y reconstruir el mensaje original.
\item Mostrar el mensaje decodificado en un display y publicarlo en una interfaz web servida por el ESP8266.
\item Transmitir código Morse mediante un buzzer, habilitando la comunicación bidireccional.
\end{itemize}
\end{alertblock}

%----------------------------------------------------------------------------------------
% INTRODUCCIÓN
%----------------------------------------------------------------------------------------
\begin{block}{1. Introducción}
En un pequeño pueblo de una zona aislada, las comunicaciones convencionales se ven afectadas por cortes de energía, falta de cobertura celular y ausencia de infraestructura de Internet. Como alternativa se propone un sistema de comunicación mediante \textbf{señales acústicas} y \textbf{código Morse}.

El ESP8266 recibe por micrófono los tonos generados por otro dispositivo (por ejemplo un teléfono), identifica puntos, rayas y pausas analizando la señal, y decodifica el mensaje.

\textbf{Restricción clave del trabajo:} la información viaja \textbf{exclusivamente por sonido}. El Wi-Fi del ESP8266 se usa \textbf{únicamente} para visualizar en el navegador el mensaje ya recibido y decodificado; no es un canal de datos.
\end{block}

%----------------------------------------------------------------------------------------
% DIAGRAMA EN BLOQUES
%----------------------------------------------------------------------------------------
\begin{figure}[h]
% Imagen necesaria: documentation/images/diagramaDeBloques.png
% \includegraphics[width=23cm, height=33cm]{diagramaDeBloques.png}
\begin{center}
\fbox{\begin{minipage}{0.90\linewidth}
\centering
\vspace{0.5cm}
\pendiente{imagen del diagrama en bloques --- guardar como
\texttt{Arq-Comp/documentation/images/diagramaDeBloques.png}}
\vspace{0.9cm}

{\small
SONIDO $\rightarrow$ MICRÓFONO $\rightarrow$ ADC $\rightarrow$ GOERTZEL\\[0.4ex]
$\rightarrow$ PUNTO / RAYA / PAUSA\\[0.4ex]
$\rightarrow$ DECODIFICADOR MORSE $\rightarrow$ MENSAJE\\[0.4ex]
$\rightarrow$ DISPLAY + INTERFAZ WEB
}
\vspace{0.6cm}
\end{minipage}}
\end{center}
\caption{Diagrama en bloques del sistema completo}
\end{figure}

\end{column}

\begin{column}{\sepwid}\end{column}

%============================================================
% COLUMNA 2 (doble ancho)
%============================================================
\begin{column}{\twocolwid}

\begin{columns}[t,totalwidth=\twocolwid]

%--------------------- Subcolumna 2.1 ---------------------
\begin{column}{\onecolwid}\vspace{-.6in}

%----------------------------------------------------------------------------------------
% MARCO TEÓRICO
%----------------------------------------------------------------------------------------
\begin{block}{2. Marco Teórico}

\subsection*{Código Morse y unidades temporales}
El código Morse codifica cada letra como una secuencia de \textbf{puntos} y \textbf{rayas}. La duración de referencia es la unidad de tiempo $u$: el punto dura $1u$, la raya $3u$; la pausa entre símbolos de una misma letra es $1u$, entre letras $3u$ y entre palabras $7u$. Como el sistema no conoce la velocidad del transmisor, en lugar de tiempos absolutos se usan \textbf{umbrales de duración en ms}, calibrables durante las pruebas.

\subsection*{Algoritmo de Goertzel}
Goertzel calcula la energía de \textbf{un solo bin} de la DFT mediante una recursión de segundo orden:
\[
\mathrm{coeff} = 2\cos\!\left(\frac{2\pi k}{N}\right),\qquad
s_n = x_n + \mathrm{coeff}\,s_{n-1} - s_{n-2}
\]
\[
P = s_{N-2}^{2} + s_{N-1}^{2} - \mathrm{coeff}\,s_{N-1}s_{N-2},\qquad
|X_k| = \sqrt{P}
\]
\textbf{Por qué alcanza y no hace falta una FFT completa:} solo interesa saber si hay energía en la frecuencia del tono. Goertzel recorre las muestras una vez, con una multiplicación-acumulación por muestra y \textbf{un único coseno} calculado fuera del lazo; no requiere el buffer complejo de $N$ puntos ni el ordenamiento de una FFT. En un ESP8266 eso se traduce en menos RAM y menos ciclos para el mismo resultado sobre el bin de interés.

\subsection*{ADC del ESP8266 y acondicionamiento}
El ADC integrado es de \textbf{10 bits} (0 a 1023) y su rango de entrada es de \textbf{0 a 1{,}0\,V}, no 0 a 5\,V. El módulo de sonido entrega la señal montada sobre un \textbf{offset de continua}, por lo que la entrada se acopla en AC y se polariza cerca de 0{,}5\,V; en el software se \textbf{resta la media del bloque} antes de Goertzel. Si el pico llega al fondo de escala, la demo avisa de saturación del ADC y hay que bajar la ganancia del potenciómetro.

\subsection*{Parámetros de muestreo elegidos}
\begin{itemize}
\item $F_s = 4000$\,Hz y $N = 400$ muestras: ventana de \textbf{100\,ms}, resolución $F_s/N = 10$\,Hz.
\item Tono objetivo de \textbf{700\,Hz} $\Rightarrow k = 70$ exacto: el tono cae justo en el bin, sin fuga espectral.
\item Compromiso de diseño: una ventana más larga mejora la resolución en frecuencia, pero un punto de $\sim$60\,ms ya no entraría completo en el bloque.
\end{itemize}

\end{block}

%----------------------------------------------------------------------------------------

\end{column}

%--------------------- Subcolumna 2.2 ---------------------
\begin{column}{\onecolwid}\vspace{-.6in}

%----------------------------------------------------------------------------------------
% METODOLOGÍA
%----------------------------------------------------------------------------------------
\begin{block}{3. Metodología}

Trabajo experimental aplicado, desarrollado en \textbf{dos etapas}: primero una demo aislada del DSP y después la integración del sistema completo.

\subsection*{Componentes}
\begin{itemize}
\item ESP8266 en modo \textbf{Access Point}: SSID \texttt{MorseESP8266}, clave \texttt{Morse123}, servidor HTTP en el puerto 80 (IP 192.168.4.1).
\item Módulo de sonido con salida \textbf{digital} (detección por umbral, potenciómetro de ganancia) y salida \textbf{analógica} para la demo de Goertzel.
\item Buzzer pasivo, pulsador o llave Morse y LED indicador.
\end{itemize}

\subsection*{Herramientas}
\begin{itemize}
\item Firmware: Arduino IDE (core ESP8266).
\item Interfaz web: Node.js con \texttt{html-minifier-terser}; \texttt{npm run build:web} minifica \texttt{client/src} y genera el header \texttt{web\_page.h}.
\end{itemize}

\subsection*{Estrategia de validación}
\begin{enumerate}
\item \textbf{Demo aislada} (\texttt{goertzel\_demo}): sin Wi-Fi ni display, solo monitor serie. Se observan las estadísticas de amplitud en silencio y con tono para fijar el umbral sobre la placa real.
\item \textbf{Integración}: los umbrales validados se trasladan al firmware principal y se verifica la cadena completa con la interfaz web.
\end{enumerate}

\subsection*{Conexiones del sistema}
{\small
\begin{tabular}{l l l}
\toprule
\textbf{Componente} & \textbf{NodeMCU} & \textbf{GPIO}\\
\midrule
Pulsador / llave Morse & D6 & 12\\
Sensor de sonido (digital) & D7 & 13\\
LED indicador & D4 & 2\\
Buzzer & D5 & 14\\
Sensor de sonido (analógico) & A0 & ADC (demo)\\
\bottomrule
\end{tabular}
}

GPIO12 se configura con \texttt{INPUT\_PULLUP}. A0 es la única entrada analógica del ESP8266 y se usa en la demo de Goertzel.

\end{block}

%----------------------------------------------------------------------------------------

\end{column}

\end{columns}

%----------------------------------------------------------------------------------------
% RESULTADOS DESTACADOS
%----------------------------------------------------------------------------------------
\begin{alertblock}{4. Resultados Destacados}

Se reporta únicamente lo implementado y verificable en el repositorio:

\begin{itemize}
\item \textbf{Decodificación A--Z funcionando:} diccionario Morse completo (26 letras), buffer del mensaje de 80 caracteres, tanto por sensor de sonido como por botón manual.
\item \textbf{Clasificación por duración:} pulso menor a 200\,ms $\rightarrow$ punto; 200\,ms o más $\rightarrow$ raya; silencio de 800\,ms o más cierra la letra; filtro antirrebote de 40\,ms en el sensor y descarte de pulsos menores a 30\,ms.
\item \textbf{Interfaz web operativa:} dos paneles (sensor y botón manual), refresco por \textit{polling} cada 300\,ms, con mensaje decodificado, código en curso, último código y estado del LED.
\item \textbf{Comunicación bidireccional:} al cerrar una letra recibida por micrófono, el ESP8266 la re-emite por buzzer (punto 100\,ms, raya 300\,ms, 2000\,Hz); el botón manual hace sonar el buzzer mientras se mantiene pulsado.
\item \textbf{Un solo origen para la UI:} \texttt{client/src} se minifica y se genera \texttt{WEB\_PAGE[] PROGMEM}, servido con \texttt{send\_P()}. Antes el \texttt{.ino} tenía una copia duplicada del HTML que ya había divergido.
\item \textbf{Demo de Goertzel con parámetros cerrados:} $F_s = 4000$\,Hz, $N = 400$, $k = 70$, ventana de 100\,ms, histéresis de 2 bloques y aviso de saturación del ADC.
\end{itemize}

\textbf{Alcance actual del DSP:} el firmware principal decodifica sobre la salida \textbf{digital} del módulo de sonido, mientras que el Goertzel trabaja sobre la salida analógica en la demo aislada. La integración del Goertzel analógico al firmware principal queda \pendiente{integrar el Goertzel del demo en el firmware principal}.

\textbf{Métricas de calidad todavía no medidas:} tasa de acierto de caracteres \pendiente{tasa de acierto de caracteres, indicar método y cantidad de ensayos}; velocidad máxima soportada \pendiente{WPM máximo soportado}; alcance acústico \pendiente{distancia máxima de recepción en metros}; latencia \pendiente{latencia de detección en ms}; tiempo de decodificación \pendiente{tiempo de decodificación por letra en ms}.

\end{alertblock}

%----------------------------------------------------------------------------------------
% MONTAJE FÍSICO (ANCHO COMPLETO)
%----------------------------------------------------------------------------------------
\begin{block}{5. Montaje Físico}
\begin{figure}[h]
% Imagen necesaria: documentation/images/montajefisico.png
% \includegraphics[width=30cm, height=14cm]{montajefisico.png}
\begin{center}
\fbox{\begin{minipage}{0.92\linewidth}
\centering
\vspace{1.8cm}
\pendiente{fotografía del montaje físico --- guardar como
\texttt{Arq-Comp/documentation/images/montajefisico.png}}
\vspace{1.8cm}
\end{minipage}}
\end{center}
\caption{Montaje físico del prototipo}
\end{figure}
\end{block}

\end{column}

\begin{column}{\sepwid}\end{column}

%============================================================
% COLUMNA 3
%============================================================
\begin{column}{\onecolwid}

%----------------------------------------------------------------------------------------
% CONCLUSIONES
%----------------------------------------------------------------------------------------
\begin{block}{6. Conclusiones}

\begin{itemize}
\item Se implementó un sistema \textbf{bidireccional} sobre ESP8266: recibe y decodifica Morse a partir de sonido, y transmite Morse con el buzzer.
\item La restricción del enunciado se cumple: el canal de datos es \textbf{acústico} y el Wi-Fi se limita a publicar el estado para visualización.
\item El AP propio (\texttt{MorseESP8266}) hace que el sistema no dependa de infraestructura de red externa.
\item La clasificación por duración con umbrales calibrables resultó suficiente para la salida digital del módulo de sonido.
\item Goertzel reduce el costo frente a una FFT completa cuando solo interesa un tono; la demo aislada permitió validar sus parámetros antes de integrar.
\item Mantener la interfaz web en un único origen evita que la página embebida en el firmware vuelva a divergir del HTML editado.
\item Robustez frente a ruido y alcance real del enlace: \pendiente{conclusión cuantitativa cuando se midan tasa de acierto y distancia}.
\end{itemize}

\end{block}

%----------------------------------------------------------------------------------------
% REFERENCIAS
%----------------------------------------------------------------------------------------
\begin{block}{7. Referencias}

\begin{itemize}
\item Arduino. \emph{Arduino Documentation}. \url{https://docs.arduino.cc}
\item ESP8266 Arduino Core. \emph{Documentation}. \url{https://arduino-esp8266.readthedocs.io}
\item Cátedra Arquitectura de Computadoras, Universidad Nacional de Hurlingham (2026). \emph{Sistema de comunicación Morse mediante audio} (enunciado del TP).
\item Goertzel, G. (1958). An Algorithm for the Evaluation of Finite Trigonometric Series. \emph{American Mathematical Monthly}, 65(1), 34--35.
\item Simulación: \pendiente{link de simulación}
\end{itemize}

% NOTA: Wokwi no soporta ESP8266. Hay que definir la herramienta de
% simulación de este TP antes de completar el marcador anterior.

\end{block}

%----------------------------------------------------------------------------------------
% CONTACTO
%----------------------------------------------------------------------------------------
\setbeamercolor{block alerted title}{fg=black,bg=norange}
\setbeamercolor{block alerted body}{fg=black,bg=white}

\begin{alertblock}{Contacto}
\begin{itemize}
\item \textbf{Emails}:
  \begin{itemize}
  \item cristianjuarezdev@gmail.com
  \item hernannicolaus@gmail.com
  \item hzambrano110@gmail.com
  \end{itemize}
\item \textbf{Teléfonos}:
  \begin{itemize}
  \item +54 9 11 3272 7298
  \item +54 9 11 3251 8092
  \item +52 55 5060 8768
  \end{itemize}
\end{itemize}
\end{alertblock}

\end{column}

\end{columns}
\end{frame}

\end{document}
